\chapter{Implementación}

En este capítulo se va a explicar cómo se ha desarrollado el sistema, qué decisiones técnicas se han tomado a lo largo 
de la implementación y qué problemas se han tenido durante el proceso.

Antes de entrar en la explicación del sistema, con las reuniones realizadas con Juan Luis Jiménez Laredo y \textbf{Pepe Agricultor ???}, se pudieron extraer los requisitos funcionales
del sistema y acotar hacia donde se iba a dirigir el desarrollo de la aplicación. Una vez fijadas las funcionalidades más importantes que iba a tener el sistema, se comenzó priorizando
y organizando el desarrollo en cada una de las iteraciones. 

Para explicar el desarrollo, se va a seguir el orden fijado por cada sprint.

\section{Sprint 1}

Para que se pueda acceder a un sistema, es necesario que la ``persona'' esté registrada en ella. Por ello, en este sprint se fijó la implementación del \textbf{registro de un usuario} y del 
\textbf{inicio de sesión}.

Primero se elaboró el esquema de la base de datos. Analizando el uso que iba a tener la aplicación por parte de los usuarios, se fijó un sistema de roles basado en dos tipos: \textbf{admin} y 
\textbf{basic} (ya explicados anteriormente).

Teniendo en cuenta esto y que un usuario tiene que estar representado también a través de una entidad, el esquema desarrollado en la primera iteración fue el sisguiente:

\textbf{INSERTAR IMAGEN DE LA PRIMERA PARTE DEL ESQUEMA}.

La entidad usuario representa al usuario como tal que accede y utiliza la aplicación. Datos como el correo electrónico o el NIF son almacenados en el sistema.

En cuanto a los roles, estos son almacenados en una tabla, en la que solo se guarda el nombre que identifica al rol.

Para poder establecer la relación entre un usuario y el rol que tiene, hay una tabla intermedia que corresponde a una ``relación''. Esta tabla almacena el ID del usuario con el ID del rol, 
estableciendo la relación de esta forma. Esta tabla intermedia permite relacionar a un usuario con varios roles, de manera que un usuario administrador tendrá rol administrador y rol básico,
mientras que un usuario normal solo tendrá rol básico.

Para crear las entidades y sus relaciones en la base de datos de PostgreSQL se ha realizado con Spring Boot directamente, ya que el framework tiene un \hyperlink{orm}{ORM} que facilita esta
implementación. Para ello, es necesario previamente establecer la conexión con la base de datos. Esto se hace a través de un archivo conocido como \textbf{application.properties}, donde se fija
el puerto, nombre de usuario, base de datos, entre otras más propiedades.

Establecidas las propiedades, se procede a crear las entidades. El proyecto ha sido organizado en distintos directorios, cuyo significado y contenido se irá explicando a lo largo de este capítulo.
Teniendo en cuenta esto, las entidades han sigo guardadas en un directorio llamado \textbf{models}, representando los modelos (entidades) del sistema. 

Una entidad en Spring Boot es creada como una clase que tiene la anotación \textbf{\textit{@Entity}}. Esto indica a Spring Boot que la clase va a ser un componente del tipo Entidad. Además de esta
anotación, se han añadido otras más como: \textbf{@Getter} (crea los métodos \textit{gets} para cada atributo), \textbf{@Setter} (crea los métodos \textit{setters} para todos los atributos) y
\textbf{@Table} (permite inidicar el nombre que tiene la entidad o tabla en la base de datos).

Siguiendo este patrón, se crearon las dos entidades junto con la relación entre ellas. Además de estas, se creó una entidad más que almacena los tokens de confirmación generados cuando un usuario se 
registra, sabiendo de esta forma qué token pertenece a cada usuario y saber si su registro está confirmado o no.

A continuación, se crearon los servicios que van a controlar el registro de usuarios y el inicio de sesión, siendo principalmente: \textbf{UserService} y \textbf{AuthService}. 

\subsection{AuthService}

Este servicio es el encargado de manejar la lógica de negocio entre el \textbf{registro} de un usuario del sistema o de un administrador de entidad, \textbf{confirmación} de un usuario registrado en el 
sistema, \textbf{inicio de sesión} de un usuario y \textbf{refresco} o renovación de los \textit{tokens} de acceso.

Entrando en más detalle, cada uno de los métodos que componen la clase tienen la siguiente funcionalidad:

\begin{itemize}
    \item \textbf{registerSystemUser:} registra un usuario en el sistema. A partir de los datos que obtiene, los envía a UserService para realizar el registro. Una vez realizado el registro, crea la URL
    para confirmar la cuenta y se envía al usuario a través del EmailService.

    \item \textbf{confirm:} a través del token que viaja en la URL del correo electrónico, se obtiene el usuario relacionado con el token. Obtenido el usuario, se habilita en el sistema permitiendo que 
    pueda acceder a la aplicación durante el inicio de sesión.

    \item \textbf{login:} permite la autenticación del usuario en el sistema. Utiliza el \textit{AuthenticationManager} de Spring Boot, que es el encargado de realizar la autenticación buscando al usuario
    en la base de datos de a cuerdo al correo electrónico y la contraseña. Si el usuario es autenticado, se obtiene a través de un objeto del tipo \textit{UserDetails}, que identifica al usuario. Con la 
    información del usuario se crean los tokens de acceso y refresco a través del JwtService. 

    \item \textbf{refresh:} permite generar nuevos tokens de acceso a la aplicación. Para ello, utiliza el token de refresco. Con este token puede extraer la información codificada y buscar si el usuario 
    sigue existiendo en el sistema y si no ha sido expirado. Si todo es correcto, se crean los nuevos tokens.
\end{itemize}

Para establecer la seguridad en la aplicación e indicar quién puede acceder a qué endpoints y cuáles requieren autenticación, se utiliza un archivo llamado \textbf{SecurityConfig}. En este archivo se fijan 
propiedades de seguridad como la configuración de \hyperlink{cors}{CORS}, si se maneja estado entra las sesiones, cómo son las excepciones que se manejan en caso de que el acceso sea denegado (falta de 
permisos o autorización) o necesidad de autenticación para ejecutar la funcionalidad. La parte más importante es la siguiente:

\begin{verbatim}
.authorizeHttpRequests(request -> request
    .requestMatchers(HttpMethod.POST, "/api/auth/register/entityAdmin")
        .authenticated()

    // Allow all requests to auth endpoints
    .requestMatchers("/api/auth/**").permitAll()
    
    .requestMatchers("/api/admin/**")
        .hasRole("ADMIN")
    .requestMatchers(HttpMethod.POST, "/api/entity/")
        .hasRole("ADMIN")
    .requestMatchers(HttpMethod.DELETE, "/api/plot/{id}")
        .hasRole("ADMIN")
    .requestMatchers("/api/entity/{entityId}/**")
        .hasAnyRole("ADMIN", "BASIC")
    .anyRequest()
    .authenticated()
)
\end{verbatim}

En esta parte, como se puede ver, se fija qué endpoints del controlador se pueden ejecutar según el rol, cuales puede ejecutar cualquier usuario y cuales requieren autenticación.

Por último y no menos importante, se pueden añadir filtros que Spring Boot necesite aplicar antes de comprobar si el usuario que ejecuta la petición HTTP tiene acceso a un recurso y si está autenticado o no.
El filtro implementado es \textbf{\textit{JwtAuthFilter}}. Este filtro se va a aplicar antes de ejecutar el contenido del endpoint del controlador (llamada la servicio), de manera que comprueba si el usuario
ha enviado en la cabecera de la petición el token, si es válido o no y si el usuario sigue existiendo en el sistema. En caso de ser correcto, autentica al usuario guardando la información en el contexto de la
petición y continuando Spring Boot el flujo del programa; en caso contrario, continua con la cadena de filtros y pasaría al contenido del archivo SecurityConfig, lanzando una excepción por falta de permisos
o autenticación.

En el siguiente diagrama se puede ver de forma gráfica el proceso de autenticación de un usuario:

\textbf{INSERTAR IMAGEN DE FLUJO DE AUTENTICACION}

\subsection{JwtService}

Este servicio es el encargado de crear los tokens de acceso a la aplicación y de refresco (\textit{refresh}). Para crear el token es necesario establecer un \textbf{tiempo de expiración}, otro tiempo de expiración para 
el de refresco y una \textbf{clave privada} con la que se firmará el token durante su creación. Esta clave es creada a partir de algoritmos como \textit{HS256} y permite verificar en el servidor que el token no ha sido manipulado durante su transmisión y que proviene de una 
fuente legítima, evitando falsificaciones y garantizando de esta manera la integridad y autenticidad.

La creación del token tiene la siguiente estructura:

\begin{verbatim}
    Jwts.builder()
        .id(user.getId().toString())
        .claims(Map.of(
            "firstname", user.getFirstname(),
            "userId", user.getId()
        ))
        .subject(user.getEmail())
        .issuedAt(new Date(System.currentTimeMillis()))
        .expiration(new Date(System.currentTimeMillis() + expiration))
        .signWith(getSignKey())
        .compact();
\end{verbatim}

\textbf{Jwts} es una clase del paquete \textit{jsonwebtoken} que permite construir tokens en Java. Lo que hacen estas líneas es lo siguiente:

\begin{enumerate}
    \item Se añade el id del usuario obtenido al realizar un inicio de sesión en la aplicación.
    \item Se introduce información adicional (\textit{claims}), como puede ser el nombre del usuario.
    \item En el campo \textit{subject} se añade el email del usuario, ya que es la forma de identificar de manera única al usuario sobre el que se crea el token.
    \item A continuación se introduce el instante de tiempo en el que se ha creado el token (\textit{issuedAt}) y en qué instante expirará el token (\textit{expiration}).
    \item Finalmente se firma el token con la clave privada generada.
\end{enumerate}

Tanto la creación del token de acceso como del token de refresco se hace utilizando el mismo código explicado. Lo único que habría que cambiar es el tiempo de expiración
de cada token, teniendo el token de refresco un tiempo de expiración mayor que el de acceso (actualmente el token de acceso tiene un tiempo de expiración de 1 hora y el
de refresco de 1 semana).

Además de la creación de tokens, también va a ser el encargado de comprobar si un token es válido o no. Para saber si un token es válido se comprueba si no ha expirado
el tiempo (el instante actual es anterior al tiempo de expiración) y que el usuario al que corresponde el token es el mismo que el que está realizando la petición (comprueba
que no se esté usando un token de otro usuario).

\subsection{EmailService}

Este servicio es el encargado de realizar el envío de un mensaje de confirmación de registro a través del correo electrónico. El envío se hace a través de una función
\textbf{asíncronca}, evitando de esta forma un bloqueo y permitiendo continuar la ejecución del código mientras el correo está siendo enviado. La función es la siguiente:

\begin{verbatim}
    try {
        MimeMessage mimeMessage = mailSender.createMimeMessage();
        MimeMessageHelper helper = new MimeMessageHelper(mimeMessage, 
                                    true, "UTF-8");

        helper.setText(emailBody(
            url, 
            firstname, 
            templatePath,
            emailText             
        ), true);
        helper.setFrom(from != null ? from : "");
        helper.setTo(to != null ? to : "");
        helper.setSubject(EMAIL_SUBJECT);
        mailSender.send(mimeMessage);
    } catch(MessagingException e) {
        throw new MailSenderException(to);
    }
\end{verbatim}

La funcionalidad es la siguiente:

\begin{enumerate}
    \item Se crea un objeto \textbf{mimeMessage} a partir de un JavaMailSender, siendo el que contendrá la estructura y propiedades
    del mensaje que se va a enviar.
    
    \item A partir de un \textbf{helper} se añadirá el cuerpo del mensaje, que será un archivo HTML junto con algunos parámetros adicionales
    como la URL que corresponde al endpoint de \textit{confirm}.

    \item Por último, antes de ser enviado el mensaje, se establece el correo del destinatario y el correo del usuario que lo envía, junto con
    un mensaje de cabecera.
\end{enumerate}

\subsection{UserService}

Este servicio es el encargado de ejecutar todas las funcionalidades relacionadas con un usuario. Algunas de estas funcionalidades son: registro
y creación de un usuario del sistema (\textit{createSystemUser}), eliminación de usuario y obtención del usuario según ID o email. 

El método de \textbf{register} es llamado desde AuthService y es el encargado de crear un usuario con el rol \textit{basic} y almacenarlo en la 
base de datos. Hay que destacar que en este método es donde se crea el \textbf{token de confirmación}, creado a través de la clase UUID que permite 
generar identificadores únicos universales, muy recomendables para claves primarias:

\begin{verbatim}
    String token = UUID.randomUUID().toString();
    LocalDateTime tokenExpiresAt = LocalDateTime.now()
        .plusHours(confirmTokenExpiresHours);
\end{verbatim}


Por otro lado, el método de \textbf{createSystemUser}, realiza lo mismo que el método de registro con la diferencia de que se 
pueden establecer los roles que tendrá el usuario (admin, basic o ambos).

\subsection{Relación entre los componentes}

Además de los servicios anteriores, también han sido implementados los siguientes componentes:

\begin{itemize}
    \item \textbf{Controladores}: estos componentes son clases con la anotación \textbf{\textit{@RestController}} de Spring Boot. Esta
    anotación indica a Spring Boot que el componente creado es un controlador que va a devolver objetos de tipo JSON. Cada controlador
    contiene distintos endpoints, siendo estos los realizan las llamadas a los servicios mencionados anteriormente. 
    
    Tomando como ejemplo el controlador \textbf{AuthController}, la estructura es la siguiente:

    \begin{verbatim}
        @RestController
        @RequestMapping(value = "/api/auth")
        @AllArgsConstructor
        public class AuthController {

            private AuthService authService;

            @PostMapping("/register")
            public ResponseEntity<UserDto> registerSystemUser(
                @Valid @RequestBody RegisterRequest request
                ) {
                return authService.registerSystemUser(request);
            }

            @PostMapping("/register/entityAdmin")
            public ResponseEntity<UserDto> registerEntityAdminUser(
                @Valid @RequestBody RegisterRequest request
                ) {
                return authService.registerEntityAdminUser(request);
            }
            
            @GetMapping("/confirm")
            public ResponseEntity<UserDto> confirm(
                @RequestParam String token
                ) {
                return authService.confirm(token);
            }

            @PostMapping("/login")
            public ResponseEntity<TokenResponse> login(
                @Valid @RequestBody LoginRequest request
                ) {
                return authService.login(request);
            }
            
            @PostMapping("/refresh")
            public ResponseEntity<TokenResponse> refresh(
                @RequestHeader(
                    HttpHeaders.AUTHORIZATION
                ) String authHeader
            ) {
                return authService.refresh(authHeader);
            }
        }
    \end{verbatim}

    Como se puede ver, aparece la anotación mencionada además de otras, como por ejemplo \textbf{\textit{@PostMapping}}, que 
    indica que ese endpoint lleva una petición HTTP de tipo ``POST''; \textbf{\textit{@Valid}}, que permite validar que los datos
    del JSON enviados son correctos, \textbf{\textit{@RequestParam}} o \textbf{\textit{@RequestBody}}, que indica que un JSON es 
    enviado en el cuerpo de la petición o que se han pasado parámetros en la URL del endpoint.

    Además del resto de controladores (UserController, EntityController, etc), hay un controlador especial que es el encargado de 
    manejar las excepciones que son lanzadas durante el sistema. El controlador ha sido implementado de la siguiente forma:

    \begin{verbatim}
        @RestControllerAdvice
        public class ErrorResponseHandler {

            @ExceptionHandler({
                UserAlreadyExistsException.class,
                UserIsEnabledException.class
            })
            public ResponseEntity<ErrorDto> userException(
                Exception ex
            ) {
                ErrorDto error = new ErrorDto(
                    "Excepción relacionada con el estado 
                    del usuario almacenado", 
                    ex.getMessage(), 
                    HttpStatus.CONFLICT.value()
                );

                return ResponseEntity.status(HttpStatus.CONFLICT)
                    .body(error);
            }

            @ExceptionHandler(DataIntegrityViolationException.class)
            public ResponseEntity<ErrorDto> integrityViolation(
                Exception ex
            ) {
                ErrorDto error = new ErrorDto(
                    "Error al ejecutar una operación SQL", 
                    ex.getMessage(), 
                    HttpStatus.CONFLICT.value()
                );

                return ResponseEntity.status(HttpStatus.CONFLICT)
                    .body(error);
            }
            
            @ExceptionHandler(HttpMessageNotReadableException.class)
            public ResponseEntity<ErrorDto> messageNotReadable(
                Exception ex
            ) {
                ErrorDto error = new ErrorDto(
                    "Error en la lectura del mensaje", 
                    ex.getMessage(), 
                    HttpStatus.CONFLICT.value()
                );

                return ResponseEntity.status(HttpStatus.CONFLICT)
                    .body(error);
            }

            @ExceptionHandler(MailSenderException.class)
            public ResponseEntity<ErrorDto> mailSenderError(
                Exception ex
            ) {
                ErrorDto error = new ErrorDto(
                    "Error durante el envío del email de confirmación", 
                    ex.getMessage(), 
                    HttpStatus.INTERNAL_SERVER_ERROR.value()
                );

                return ResponseEntity
                    .status(HttpStatus.INTERNAL_SERVER_ERROR)
                    .body(error);
            }

            @ExceptionHandler({
                ConfirmTokenNotExistsException.class,
                RoleNotExistsException.class,
                UserNotFoundException.class,
                EntityNotFoundException.class
            })
            public ResponseEntity<ErrorDto> entityNotExists(
                Exception ex
            ) {
                ErrorDto error = new ErrorDto(
                    "La entidad no existe en la base de datos", 
                    ex.getMessage(), 
                    HttpStatus.NOT_FOUND.value()
                );

                return ResponseEntity.status(HttpStatus.NOT_FOUND)
                    .body(error);
            }

            ...
        }
    \end{verbatim}

    Este controlador le indica a Spring Boot que es el encargado de manejar las excepciones a través de la anotación 
    \textbf{\textit{@RestControllerAdvice}}. Con la anotación \textbf{\textit{@ExceptionHandler}} se va indicar qué
    tipo de excepción va a tratar el método y que respuesta va a devolver. Las respuestas tienen la misma estructura 
    a través de un DTO (explicado a continuación) llamado ``ErrorDto''.

    \item \textbf{Repositorios}:
    \item \textbf{DTOs}:
\end{itemize}

Los flujos de comunicación o como interactaún entre ellos es de la siguiente forma:




\section{Sprint 2}

\section{Sprint 3}

\section{Sprint 4}





Hablar en algún momento de los PreAuthorize

